Dalam perancangan dan pengembangan aplikasi pengarsipan data otomatis serta portal \textit{on-demand retrieval} ini, terdapat beberapa batasan masalah yang ditetapkan sebagai fokus penelitian, antara lain:
\begin{enumerate}
    \item Rentang waktu dan kriteria pengarsipan data historis mahasiswa---baik berdasarkan Tahun Akademik, Nomor Induk Mahasiswa (NIM), maupun parameter akademik lainnya---bersifat dinamis dan fleksibel sesuai dengan konfigurasi sistem yang dapat diatur oleh administrator. Kriteria pengarsipan tidak terikat pada aturan baku jangka waktu tertentu, melainkan dapat disesuaikan dengan kebijakan operasional Pusdatin UNTAR yang berlaku pada saat implementasi.

    \item Eksekusi pemindahan data tidak dilakukan secara langsung pada server utama Pusdatin UNTAR melalui eksekusi lokal, melainkan menggunakan arsitektur \textit{Client-Worker} terdistribusi yang berkomunikasi melalui layanan HTTP API JSON (\textit{agent.ashx}). Mekanisme ini memisahkan tanggung jawab inisiasi tugas (\textit{client}) dari eksekusi pemindahan data aktual (\textit{worker}) secara terstruktur.

    \item Eksekusi antrean tugas (\textit{task queue}) pengarsipan dijalankan menggunakan SQL Server Agent dan skrip PowerShell yang dipasang pada mesin terpisah (laptop eksternal / \textit{external machine}) yang bertindak sebagai \textit{primary agent}. Pendekatan ini dipilih secara khusus untuk mengatasi batasan infrastruktur dan kapasitas komputasi pada server kampus Pusdatin UNTAR yang tidak memungkinkan penggunaan mekanisme penjadwalan lokal secara langsung.

    \item Portal \textit{web} pengarsipan dan pencarian data dikembangkan menggunakan kerangka kerja ASP.NET (C\# / Web Forms) dan terintegrasi dengan basis data Microsoft SQL Server sebagai sistem manajemen basis data utama.

    \item Seluruh proses pengujian fungsionalitas dan validasi sistem aplikasi---mencakup pengujian alur pemindahan data, pengujian antrean tugas, serta pengujian pencarian \textit{on-demand}---sepenuhnya dieksekusi di dalam lingkungan pengembangan dan uji coba (\textit{development/staging environment}), bukan pada basis data produksi Pusdatin UNTAR yang sedang berjalan secara langsung.
\end{enumerate}
